\section{Complexity in the full binary input length}
\label{sec:complexity}

We first bound the recursion independently of the cost of solving an
ordinary mean-payoff game.  This separates the combinatorial reduction
from the numerical subroutine and shows exactly where the full binary
encoding enters the running time.

\begin{proposition}[Cost of the reduction]
\label{prop:reduction-cost}
For an arena with $n\ge1$ vertices and complete binary input length $L$,
the call $\Solve(G,(n,n))$ makes
\[
 2^{O((\log_2(n+2))^2)}
\]
calls to the ordinary mean-payoff winning-set subroutine.  Each such
call has an explicit input of length $(L+2)^{O(1)}$.  Apart from those
subroutine calls, the computation takes
\[
 2^{O((\log_2(n+2))^2)}(L+2)^{O(1)}
\]
bit operations.  All constants are absolute.
\end{proposition}

\begin{proof}
Consider the tree of calls to $\Solve$, counting its immediate base
cases as nodes.  Every child omits the minimum priority of its parent's
input arena.  Thus a path has at most $n$ edges, independently of the
numerical values of the priorities.

Each node has at most one child with both size bounds unchanged.
Every other child has one bound halved, with rounding down.  In either
phase, a child is followed by a nonempty deletion unless that phase
returns.  Each phase therefore makes at most $n+1$ child calls, giving
at most
\[
 M=2(n+1)
\]
children with a halved bound at a node.  Along a path, either one of the
two size bounds can be halved at most $1+\lfloor\log_2 n\rfloor$
times before reaching zero, where the recursion stops.  Consequently,
there are at most
\[
 h=2(1+\lfloor\log_2 n\rfloor)
\]
halving edges on any path.

The resulting path count is the same combinatorial saving that
underlies reduced-precision parity recursion~\cite{Parys2019}.
For a path of length $\ell$ with $j$ halving edges, choose their
positions and, at each one, the child index among at most $M$ halving
children.  At every remaining position there is at most one available
unchanged-bounds child.  This gives at most $\binom{\ell}{j}M^j$ paths
of that type.  As each node is specified by its root-to-node path, the
total number of calls $N_{\rm rec}$ satisfies
\begin{align}
\label{eq:recursion-count}
 N_{\rm rec}
 &\le \sum_{\ell=0}^{n}\sum_{j=0}^{\min(\ell,h)}
          \binom{\ell}{j}M^j\\
 &\le (n+1)(h+1)(nM)^h
  =2^{O((\log_2(n+2))^2)}.\nonumber
\end{align}
The second inequality uses $\binom{\ell}{j}\le n^j$ and $nM\ge1$;
it also covers $n=1$.

We next count the nonrecursive work at a node.  A phase iteration that
continues deletes at least one vertex.  The two phases together
therefore make at most $n+2$ iterations, allowing one final return
iteration for each phase.  In particular, they make $O(n+1)$ ordinary
mean-payoff calls.  Attractor computations, construction of induced
subarenas, scans for minimum priorities, and updates of the two size
bounds all take polynomially many bit operations in $L+2$.  No step
iterates through the integers between two successive priority values:
minimum priorities and equality tests are computed directly from their
binary encodings.

All mean-payoff instances use induced subarenas of the original graph,
with the weights $w(v,u)=r(v)$.  Renumbering retained vertices uses only
polynomial space and time.  Even if the reward $r(v)$ is copied
explicitly onto every retained outgoing edge, the total number of
weight bits is at most the number of edges times the largest original
reward bit length, hence polynomial in $L+2$.  The rest of each
subroutine input is also polynomial in $L+2$.  This proves the claimed
input-length bound without a dependence on the numerical magnitude
of a reward.  Combining these per-node bounds with
\eqref{eq:recursion-count} proves the proposition.
\end{proof}

\begin{proof}[Proof of Theorem~\ref{thm:main}]
The empty arena has an empty winning set and is handled immediately.
For a nonempty arena, the output of $\Solve(G,(n,n))$ is exact by
Theorem~\ref{thm:labels} and the winning-region argument at the end of
Section~\ref{sec:algorithm}.

Proposition~\ref{prop:mp-interface} implements each ordinary mean-payoff
call of input length $L'$ in
$2^{O((\log_2(L'+2))^2)}$ bit operations, uniformly and deterministically.
By Proposition~\ref{prop:reduction-cost}, there is an absolute constant
$a$ such that $L'\le(L+2)^a$ for every call, after increasing $a$ if
necessary.  Its cost is therefore
$2^{O((\log_2(L+2))^2)}$.  The number of calls and all remaining work
have the same bound, since $n$ is polynomially bounded by the explicit
input length.  Their product is again
\[
 2^{O((\log_2(L+2))^2)}.
\]
All graph operations and recursion choices admit deterministic
implementations, so the complete procedure has the stated uniform
deterministic bit complexity.  In particular, the constants do not
depend on the number of priorities or on the numerical sizes of the
binary rewards and priorities.
\end{proof}
